% Moved from the main text in the scope cut (2026-09-30); text unchanged except cross-references.
\section*{S18. Cohorts and further results in detail}\phantomsection\label{supp:S18}
\subsection*{Cohorts in detail}
\paragraph{Controlled dermoscopy cohort} ISIC 2018 Task 1/2 images with manual lesion masks and melanoma labels
from the union of the ISIC 2016 and 2017 tables (three images with conflicting labels dropped); perceptual-hash
(Hamming $\le 8$) and lesion grouping; 2,437 images for which a $65\times19$\,px ruler (0.46\% of a
$518\times518$ image) can be placed at 0/25/50/75/100\% overlap with the lesion within 10 percentage points
(train/val/test 1,849/371/371 before common support).

\paragraph{Real hair cohort (ISIC 2019)} All 25,331 ISIC 2019 training images; hair/ruler masks from the archive of
\citet{wegley2026artifacts}; lesion masks manual for HAM10000 \citep{tschandl2018ham} and predicted elsewhere by a
U-Net trained on ISIC 2018 Task 1. Two traps share one hair-free group ($\le 30$ hair pixels at native resolution):
Trap~A contains hair with overlap $r=|\text{hair}\cap\text{lesion}|/|\text{hair}|\ge0.5$, Trap~B $r<0.1$. Cells are
subsampled to the hair-free source mix (HAM/BCN/MSK) within each label, giving Trap~A 704/157/2,738/989 and Trap~B
704/157/4,327/930 images (hair-free benign/melanoma, hair benign/melanoma); five seeds of five-fold group-safe
cross-validation (leakage groups: lesion identifiers and pHash); 172 reversed-test melanomas per trap and seed. The
30-pixel threshold is sensitive: independent labels show that it marks hair on many images that experts call hair-free
(Robustness paragraphs below).

\paragraph{Thyroid ultrasound} TN3K images \citep{gong2021tn3k} with TNCD benign/malignant labels from cytological
biopsy \citep{gong2022acl} (3,493 images of 2,421 patients; 1,210 malignant) and manual nodule masks; the official
train/test split (2,879/614 images) is patient-disjoint \citep{gong2022acl}. Sonographer calipers were segmented by a frozen rule-based detector (dotted measurement lines
and `+' ends; audits below and in S4). Trap~A: calipers with $r\ge0.5$ (1,150 images); Trap~B: $r<0.1$ (326);
caliper-free: 1,801. Patient identifiers are not published, so within the pooled images leakage groups are pHash
chains (stricter embedding-based groups below).

\paragraph{Ovarian tumour ultrasound} MMOTU OTU\_2d \citep{zhao2022mmotu} (CC BY 4.0): 1,202 images with tumour
masks, labelled benign cystic (chocolate cyst, serous cystadenoma, simple cyst) versus lesions with solid
components (teratoma, theca-cell tumour, mucinous cystadenoma, high-grade serous carcinoma). The thyroid detector is
applied unchanged: Trap~A 341/276, Trap~B 91/103, caliper-free 170/178 (benign-cystic/solid).

\paragraph{Capsule endoscopy} SEE-AI frames \citep{seeai2023} (CC BY 4.0) containing only erosions (3,550) or only
polyp-like lesions (1,931); the ROI is the union of the expert lesion boxes. Contamination (debris, bubbles, bile)
is segmented by a linear probe on DINOv2 patch tokens trained on 2,160 expert-masked frames (held-out pixel accuracy
0.92, IoU 0.62; expert masks are used where available). Debris-free: $<3\%$ of the frame; Trap~A: debris $\ge10\%$
with $r\ge0.5$; Trap~B: $r<0.1$. Leakage groups: blocks of 100 consecutive frame numbers plus pHash duplicates.

\paragraph{Chest radiography} NIH ChestX-ray14 \citep{wang2017chestxray8}; lung ROIs from the TorchXRayVision
PSPNet \citep{cohen2022xrv}. \emph{Controlled cohort:} 1,906 pneumothorax images without a chest drain (NEATX
annotations) and 5,718 \emph{No Finding} images from other patients; a $100\times12$\,px radio-opaque tube is
placed at controlled overlap with the lungs (7,610 images in common support). \emph{Real drains} (in-ROI only):
29,687 images; drains are labelled by NEATX for pneumothorax-positive images and by a drain detector (precision
0.95) for negatives; environments are built within view position (AP/PA); five patient-grouped folds.

\subsection*{Further results}
\paragraph{Masking implementation and the real overlap (pre-registered)} Mean-colour fill is one way to mask. Black fill, a blurred background (16-px Gaussian), a crop to the lesion box and a
crop of the masked image all give a positive crossover in all four cohorts (16 of 16, Holm $p\le0.010$), from
$+0.091$ [$+0.064, +0.117$] for the dermoscopy box crop to $+0.367$ [$+0.318, +0.416$] for black fill on capsule
frames; implementations that keep part of the out-of-ROI content give weaker effects where that content still carries
the artifact (blur for debris and ovarian calipers, the box crop for hair and debris). The two traps are the ends of a
continuous dependence: with all artifact-bearing images binned by their real overlap $r$, the mask gain falls with $r$ in
every cohort (slopes: hair $-0.235$ [$-0.276, -0.194$], thyroid $-0.250$ [$-0.303, -0.197$], capsule $-0.489$
[$-0.539, -0.438$], ovary $-0.205$ [$-0.300, -0.109$]), and for hair it turns from a gain ($+0.124$ [$+0.085, +0.162$]
at $r<0.1$) into a harm ($-0.086$ [$-0.119, -0.056$] at $r\ge0.75$), crossing zero near $r=0.40$ (S16).

\paragraph{Replication across encoders, training regimes and foundation models} The crossover replicates with MedSigLIP (thyroid $+0.195$ [$+0.153, +0.238$], capsule $+0.335$ [$+0.286, +0.383$],
ovary $+0.178$ [$+0.109, +0.248$]), a CNN (ConvNeXt-Base: thyroid $+0.275$ [$+0.231, +0.320$], capsule $+0.320$ [$+0.268, +0.372$]) and
DINOv2 at three sizes (positive in all nine cohort--size combinations; S7), and with end-to-end fine-tuning (crossover of a fine-tuned ResNet-50: hair $+0.143$ [$+0.054, +0.232$]; capsule $+0.606$ [$+0.495, +0.716$]; thyroid ViT-S $+0.122$ [$+0.039, +0.211$], ResNet-50 $+0.063$ [$-0.022, +0.144$]). The fine-tuned ovary network did not learn the task (ERM clean AUROC 0.60 and 0.50 in Trap~A and B, 15 clusters), so it can neither show nor refute a shortcut effect and ovary is not evaluable for this check; the replication holds in all three evaluable cohorts%
.
Without any task training, the zero-shot scores of vision--language models are already swayed by artifacts, and
masking changes them as the law predicts (MedSigLIP crossovers: thyroid $+0.103$ [$+0.064, +0.143$], capsule $+0.045$
[$+0.001, +0.090$], ovary $+0.181$ [$+0.103, +0.261$]; DermLIP on hair $+0.040$ [$+0.021, +0.058$]); MedSigLIP's
zero-shot thyroid diagnosis itself is at chance (clean AUROC 0.45), so we use these scores only to show that the
effect is a property of the representation.

\paragraph{The one encoder without a hair gap} With the dermatology foundation model DermLIP and a trained linear
head, masking harms at \emph{both} hair locations (Trap~A $-0.104$ [$-0.143, -0.067$], Trap~B $-0.090$
[$-0.135, -0.046$]) and the hair crossover vanishes ($+0.014$ [$-0.037, +0.063$]), although the controlled ruler sweep
reverses with DermLIP as with DINOv2 ($+0.161$ at no overlap, $-0.263$ to $-0.309$ at 25--100\%). Two measurements
explain it. First, with DermLIP masking costs disease signal wherever the hair lies (clean AUROC $-0.066$
[$-0.089, -0.043$] in Trap~A and $-0.063$ [$-0.094, -0.033$] in Trap~B): this encoder, trained on whole dermatology
images, loses more from a masked background than DINOv2. Second, the pixel part of its hair shortcut is small in both
traps (oracle hair inpainting $+0.033$ [$+0.025, +0.042$] and $+0.040$ [$+0.015, +0.064$]), whereas balancing, which
also removes correlates, gains $+0.146$ and $+0.110$: masking removes little of the shortcut at either location, so
there is no gap to find. The linear-Gaussian account fitted to DermLIP's clean and correlated AUROC anticipates this
(predicted crossover $+0.013$, observed $+0.014$; main-text theory section). The law describes what masking can remove; when the shortcut is carried by
correlates and the encoder is fragile to masking, masking removes little at either location.

\paragraph{Operating points: registration status and capsule} \emph{Registration status for this frozen model: a pre-registered test inside a
post hoc subgroup.} The subgroup (the 78 malignant nodules of the official test split that carry an in-ROI caliper; the
same 78 in every seed) was defined after the first unaltered-data analysis, where its AUROC harm was found; the
operating-point test in it was registered before the operating points were computed, and for the fine-tuned model below
the whole test was registered before any result. The pre-registered decision rule
(loss at the maximal-balanced-accuracy threshold and at $\ge3$ of the four other operating points) is met. On capsule
endoscopy, where masking improves AUROC overall, it lowers the sensitivity for erosions carrying debris on the lesion
at the two high-sensitivity operating points (S9).

\paragraph{A fine-tuned model near the benchmark (pre-registered)} To test whether these are properties of a weak
model, we compared twelve fine-tuning recipes on validation images only, committed the choice (ConvNeXt-T pretrained on
ImageNet-22k, learning rate $10^{-4}$, 8 epochs, 224\,px) and then trained it with five seeds. Its test AUROC is 0.769,
0.004 below the lowest published cross-entropy baseline on this split (0.773; \citealp{gong2022acl}) and above the
frozen probe (0.735). All four registered tests are supported (Holm): masking lowers the AUROC on shortcut-conflicting
pairs by $-0.102$ [$-0.178, -0.026$]; it lowers sensitivity at all four of OP2--OP5; at OP2 it lowers sensitivity among
the 78 malignant nodules with an in-ROI caliper from 0.744 to 0.223 ($-0.521$ [$-0.642, -0.382$]), a subgroup registered
before any result of this model existed; and its trap crossover is $+0.187$ [$+0.111, +0.268$]. Group balancing recovers
the conflicting pairs ($+0.088$ [$+0.019, +0.158$] over masking).

\paragraph{Is the loss a threshold artefact? (pre-registered)} Re-tuning each model's threshold on the test set
itself---the best that recalibration at a deployment site can achieve---makes the overall sensitivity of the two models
equal at a test specificity of 0.80 ($-0.017$ [$-0.124, +0.070$] for the fine-tuned model), but the conflicting subgroup
still loses: from 0.544 to 0.390 ($-0.154$ [$-0.302, -0.023$]) for the fine-tuned model and from 0.505 to 0.331
($-0.174$ [$-0.313, -0.012$]) for the frozen probe (Holm $p=0.019$ each), while the malignant nodules whose caliper
points the right way tend to gain. About 30\% of the loss at validation-fixed thresholds therefore cannot be removed by
any threshold; the rest follows from transporting thresholds from validation to test data, as deployment does. At a
matched specificity of 0.90 and a matched sensitivity of 0.80 the direction is the same, with intervals that include
zero (S16).

\paragraph{New patients at natural prevalence (pre-registered)} Trained on all ISIC 2019 images and tested on ISIC 2020
without any resampling, masking lowers the AUROC on shortcut-conflicting pairs: $-0.026$ [$-0.052, -0.000$] under the
registered near-duplicate rule, which proved too loose (it removed most look-alike images of other patients), and
$-0.037$ [$-0.061, -0.012$] under a label-free calibrated rule registered as a sensitivity analysis after that result
(32,952 images). The overall AUROC is nearly unchanged ($+0.016$ [$-0.000, +0.033$]), and group balancing recovers the
conflicting pairs ($+0.072$ [$+0.048, +0.097$] over masking). No sensitivity loss at
a fixed operating point was detectable, and the cleanest in-lesion comparison (53 hair-free melanomas) was
inconclusive ($-0.005$ [$-0.057, +0.048$]): the external harm combines the removal of hair beside melanomas with the
retention of hair on benign lesions (S16).

\paragraph{Calibration} For the held-out BCN hospital masking worsens both Brier score and expected calibration error
(Table~\ref{tab:natural_clinical}); for HAM, MSK and capsule it improves every metric, in line with the AUROC results.
\begin{table}[t]\centering\scriptsize
\caption{Clinical secondary metrics on the unaltered test sets (Table~\ref{m:tab:natural}): mean (SD) over training seeds. Sensitivity and specificity at the operating point fixed on clean validation data; ECE with 10 equal-width bins; AUPRC baseline = prevalence.}\label{tab:natural_clinical}
\begin{tabular}{llcccccc}\toprule
Test set & Arm & AUROC & AUPRC & Brier & ECE & Sens & Spec \\\midrule
Thyroid (official split) ($\pi=0.38$) & ERM & 0.736 (0.009) & 0.606 (0.012) & 0.207 (0.005) & 0.082 (0.008) & 0.734 (0.047) & 0.619 (0.061) \\
 & mask & 0.731 (0.006) & 0.632 (0.007) & 0.214 (0.006) & 0.120 (0.026) & 0.347 (0.095) & 0.888 (0.041) \\
 & U-MtE & 0.742 (0.003) & 0.638 (0.007) & 0.203 (0.001) & 0.085 (0.007) & 0.425 (0.088) & 0.854 (0.050) \\
 & balanced & 0.716 (0.007) & 0.584 (0.012) & 0.219 (0.004) & 0.110 (0.007) & 0.714 (0.052) & 0.599 (0.061) \\
 & U-MtE$_{\text{bal}}$ & 0.731 (0.004) & 0.630 (0.007) & 0.207 (0.001) & 0.084 (0.006) & 0.399 (0.069) & 0.866 (0.034) \\
\midrule
ISIC, BCN held out ($\pi=0.23$) & ERM & 0.786 (0.003) & 0.586 (0.003) & 0.239 (0.006) & 0.279 (0.011) & 0.842 (0.038) & 0.519 (0.075) \\
 & mask & 0.769 (0.004) & 0.571 (0.009) & 0.286 (0.007) & 0.332 (0.009) & 0.860 (0.013) & 0.439 (0.023) \\
 & U-MtE & 0.773 (0.004) & 0.574 (0.008) & 0.280 (0.009) & 0.328 (0.012) & 0.858 (0.039) & 0.448 (0.089) \\
 & balanced & 0.781 (0.005) & 0.585 (0.009) & 0.228 (0.006) & 0.257 (0.011) & 0.811 (0.046) & 0.554 (0.087) \\
 & U-MtE$_{\text{bal}}$ & 0.774 (0.004) & 0.568 (0.008) & 0.279 (0.011) & 0.325 (0.016) & 0.860 (0.027) & 0.447 (0.051) \\
\midrule
ISIC, MSK held out ($\pi=0.19$) & ERM & 0.770 (0.008) & 0.477 (0.012) & 0.172 (0.007) & 0.161 (0.019) & 0.608 (0.049) & 0.786 (0.039) \\
 & mask & 0.790 (0.004) & 0.512 (0.009) & 0.146 (0.003) & 0.115 (0.009) & 0.593 (0.052) & 0.824 (0.037) \\
 & U-MtE & 0.783 (0.004) & 0.500 (0.008) & 0.150 (0.003) & 0.126 (0.009) & 0.566 (0.073) & 0.831 (0.049) \\
 & balanced & 0.766 (0.008) & 0.462 (0.014) & 0.178 (0.008) & 0.170 (0.020) & 0.626 (0.047) & 0.755 (0.046) \\
 & U-MtE$_{\text{bal}}$ & 0.774 (0.006) & 0.482 (0.009) & 0.154 (0.004) & 0.131 (0.010) & 0.594 (0.061) & 0.798 (0.051) \\
\midrule
ISIC, HAM held out ($\pi=0.11$) & ERM & 0.750 (0.005) & 0.353 (0.005) & 0.162 (0.006) & 0.234 (0.011) & 0.535 (0.092) & 0.813 (0.068) \\
 & mask & 0.789 (0.002) & 0.404 (0.006) & 0.143 (0.005) & 0.217 (0.011) & 0.559 (0.112) & 0.842 (0.063) \\
 & U-MtE & 0.799 (0.008) & 0.406 (0.009) & 0.145 (0.007) & 0.226 (0.016) & 0.573 (0.067) & 0.846 (0.041) \\
 & balanced & 0.747 (0.017) & 0.346 (0.019) & 0.156 (0.014) & 0.228 (0.015) & 0.539 (0.078) & 0.803 (0.076) \\
 & U-MtE$_{\text{bal}}$ & 0.802 (0.005) & 0.401 (0.009) & 0.139 (0.006) & 0.215 (0.021) & 0.606 (0.054) & 0.825 (0.037) \\
\midrule
Capsule (held-out frames) ($\pi=0.56$) & ERM & 0.889 (0.019) & 0.897 (0.046) & 0.130 (0.016) & 0.042 (0.011) & 0.770 (0.062) & 0.844 (0.055) \\
 & mask & 0.967 (0.011) & 0.971 (0.014) & 0.069 (0.013) & 0.029 (0.011) & 0.859 (0.032) & 0.933 (0.025) \\
 & U-MtE & 0.962 (0.009) & 0.967 (0.013) & 0.081 (0.017) & 0.065 (0.045) & 0.871 (0.042) & 0.897 (0.043) \\
 & balanced & 0.881 (0.024) & 0.889 (0.053) & 0.136 (0.017) & 0.051 (0.009) & 0.758 (0.087) & 0.838 (0.059) \\
 & U-MtE$_{\text{bal}}$ & 0.959 (0.008) & 0.964 (0.012) & 0.079 (0.007) & 0.044 (0.012) & 0.866 (0.048) & 0.897 (0.052) \\
\bottomrule
\end{tabular}\end{table}


\paragraph{Robustness to artifact-label definitions}
Artifact masks for calipers come from a detector (visual audit: 95\% of thyroid detections lie on real calipers;
92--97\% of caliper-free images contain no visible caliper). Re-running the thyroid and ovary analyses with
large-marker-only ($\ge 50$ px) and strict-location ($r\ge0.7$ / $r<0.05$) definitions leaves the crossover
(thyroid $+0.215$ to $+0.262$; ovary $+0.170$ to $+0.195$) and the U-MtE$_{\text{protect}}$ gain over masking
(thyroid $+0.217$ to $+0.282$; ovary $+0.043$ to $+0.067$) significant in every variant; both grow under stricter
definitions, as expected if detector errors dilute rather than create the effects.

\paragraph{Artifact labels checked without a clinician (pre-registered)}
No clinician was available to audit the labels, so we compared them with every independent source we could find, with
coverage gates and without the images on which our own labellers were trained (S16). \emph{Dermoscopy.}
Where hair is present our labels place it correctly: our lesion masks agree with the IMA++ multi-annotator masks (241
images; Dice 0.851 [0.830, 0.870]; trap class agreement 84\%), and against independent hair masks (485 images) the hair
lies on the other side of the lesion boundary in 1 of 264 trap images. But at our threshold ($>30$ pixels) the published
hair masks mark hair on 93\% of the 8,864 images that the expert-reviewed DermArtifactDB labels call hair-free
($\kappa=0.076$), more often for melanomas. Removing every image an independent source contradicts (42\% of each trap)
raises the hair crossover from $+0.147$ [$+0.111, +0.185$] to $+0.199$ [$+0.164, +0.233$]: the label noise diluted
the effect. \emph{Calipers.} Two independent labellers---an ultrasound annotation detector (BUSClean) and a medical
vision--language model (MedGemma, run locally with fixed prompts and never shown our masks)---both pass a pre-registered
informativeness rule on thyroid ($\kappa\ge0.40$) but err in opposite directions, so the registered rule that removes an
image whenever either disagrees left too few out-of-ROI caliper images to re-run the thyroid traps. Where the two agree
with each other (1,865 of 3,277 images) they agree with our cell for 94\%, and without the images both contradict the
thyroid crossover is $+0.257$ [$+0.209, +0.304$] (exploratory; $+0.239$ on our labels); the subgroup loss of
the main-text section on unaltered data holds in the 36 nodules whose in-ROI caliper both confirm. On ovarian screens the annotation
detector fires on 90\% of caliper-free images and was declared invalid, so the ovary and capsule labels could not be
checked with public data.

\paragraph{Stricter leakage grouping (pre-registered)}
TN3K/TNCD, MMOTU and SEE-AI publish no patient or video identifiers, so the main analyses group near-duplicates by
perceptual hash (and, for capsule, blocks of 100 consecutive frames). As a pre-registered sensitivity analysis we
united these groups with connected components of mean-centred DINOv2 embeddings at the most aggressive cosine
threshold that keeps every group below 5\% of the cohort (thyroid: 3,455 $\to$ 3,081 groups, $\tau=0.80$; ovary:
1,059 $\to$ 969, $\tau=0.80$; capsule: 137 $\to$ 135, $\tau=0.93$) and re-ran all primary contrasts. Eleven of twelve
keep their sign and significance: the location crossover is $+0.235$ [$+0.193, +0.275$] (thyroid), $+0.157$
[$+0.053, +0.255$] (ovary) and $+0.353$ [$+0.296, +0.410$] (capsule), and protected U-MtE still beats masking on
thyroid ($+0.251$ [$+0.215, +0.284$]) and ovary ($+0.051$ [$+0.020, +0.079$]). The exception is the small ovary
advantage of erasure over augmentation, which loses significance ($+0.024$ [$-0.014, +0.062$]). Embedding grouping
removes visually near-identical images of one patient or session; it cannot exclude differently looking images of
the same patient.

\paragraph{Rebuild on new hardware (pre-registered)}
All data were downloaded again and every derived object---caches, perceptual-hash groups, the retrained lesion U-Net
(held-out Dice 0.886; HAM10000 agreement 0.895), the caliper scan and the debris probe (IoU 0.625)---was rebuilt on a
different GPU generation (pre-registered). Trap membership changes slightly where it depends on a retrained model
(hair traps: Trap~A 2,655/942 artifact-bearing benign/melanoma images instead of 2,738/989); the thyroid and ovary
cells are identical. Every primary crossover reproduces within the pre-registered tolerance: dermoscopic hair $+0.147$ [$+0.119, +0.177$] (archived +0.145); thyroid $+0.239$ [$+0.212, +0.264$] (archived +0.230); ovary $+0.170$ [$+0.116, +0.223$] (archived +0.170); capsule $+0.377$ [$+0.339, +0.412$] (archived +0.368).%


\paragraph{Limitations in detail}
(i) \emph{Artifact labels.} Calipers are segmented by a rule-based detector and capsule debris by a patch-token probe
(IoU 0.62 against expert masks); hair masks are published semi-automatic masks. No clinician audited them. Independent
sources confirm where hair lies, but our hair-presence threshold over-calls hair, with errors that differ by diagnosis;
the thyroid caliper labels agree with two independent labellers where those agree with each other, but the registered
cleaning rule could not be tested for thyroid, and the ovary and capsule labels remain unchecked (Robustness paragraphs above).
A blinded rating of 240 images by the first author, a non-clinician, is registered and not yet done. Label errors that
do not depend on the diagnosis dilute the constructed correlation, and removing contradicted images or using stricter
label definitions strengthens rather than weakens the results (Robustness paragraphs above); the transplant experiment does not depend on the artifact-free group being
perfectly clean because the transplanted instance is identical at both locations.
(ii) \emph{Patient identifiers.} TN3K/TNCD, MMOTU and SEE-AI publish no patient or video identifiers for the pooled
images, so the trap splits group near-duplicates by perceptual hash and frame blocks (a stricter embedding grouping
leaves 11 of 12 primary estimates unchanged); the official TN3K/TNCD test split used for the unaltered-data analysis is
patient-disjoint. The one open cohort we found with patient identifiers, pathology labels and real calipers
(ThyUS2Path; 842 patients) has no nodule masks and uses a scanner interface on which our frozen caliper detector
fires on interface elements; external validation on it requires a transferred nodule segmentation and a re-validated
detector and is left for future work. For dermoscopy, external validation on new patients was possible and was run under
a pre-registered cell-count gate (ISIC 2020, patient-disjoint folds, hair masks from a segmenter trained only on
ISIC 2019): the crossover replicates ($+0.250$ [$+0.174, +0.327$]), and at natural prevalence masking lowers the AUROC on
conflicting pairs (main text, unaltered test data). External validation is limited to dermoscopy, and on ISIC 2020 no
operating-point effect was detectable.
(iii) \emph{Constructed shifts.} Reversed environments are constructed by subsampling; real deployment shifts are
less extreme, which is why we report the unaltered test distributions separately.
(iv) \emph{Absolute performance.} Frozen encoders with linear heads are below the best published fine-tuned models on
the one cohort with a directly comparable benchmark (official TNCD split: 0.735 AUROC versus 0.773--0.784 for
fine-tuned CNNs with cross-entropy and up to 0.799 with curriculum learning, \citealp{gong2022acl}). A fine-tuned
ConvNeXt-T reaches 0.769, near but below the lowest published baseline, and shows every registered harm
(main text, unaltered test data). The law replicates with end-to-end fine-tuned networks in three of four cohorts; on ovarian ultrasound the fine-tuned ResNet-50
barely learns the task and shows no significant crossover (Replication paragraph above). The other cohorts' binary
tasks have no published benchmark on the same split (Supplement~S10).
(v) \emph{Scope.} Chest radiography is a boundary case: the radiograph encoder ignored extra-pulmonary tubes, the
crossover that a second encoder shows rests on a device contrast confounded with other tubes, and real drains carry the
shortcut through correlates. The linear-Gaussian account is fitted per model and cannot anticipate a shortcut
carried by correlates that the clean and correlated environments do not expose.
(vi) \emph{Uncertainty and balance.} The original per-seed bootstrap understated CIs where seeds share test
images. All model-comparison intervals were therefore recomputed with the crossed bootstrap from regenerated runs
(main-text protocol section);
claims whose interval then included zero were rewritten, and a few secondary analyses from runs that could not be
regenerated are given as point estimates only. Matching left within-label standardised
differences of up to 0.19 for some covariates. (vii) \emph{Transplant.} Pasted patterns, including neutral tissue,
become learnable shortcuts; the transplant establishes location dependence on fixed images, not the size of real
artifacts' effects in their natural context (for hair, whose natural shortcut is carried mostly by correlates, the
difference is large). (viii) \emph{Pre-registration.} Every registration is committed before its outcome results (registry with commit hashes
and times in Supplement~S1); commit times are set by the committing machine, so only the registrations that were
pushed to GitHub before their runs (the second review round and every later round) carry an external time stamp, and
the 16-test primary family was assembled after the individual results. Two later analyses are sensitivity analyses
registered after the result they examine (the calibrated duplicate rule, the cleaner hair groups), and a few are
exploratory (the consensus caliper labels). (ix) \emph{Operating points.} Most of the thyroid sensitivity loss at
validation-fixed thresholds is threshold transport; the part that survives re-tuning has a wide interval
($-0.154$ [$-0.302, -0.023$]). (x) \emph{Remedies.} No remedy dominates masking: in two registered rounds the best candidates fail on one cell, ISIC
2020 overall, by 0.01--0.02 AUROC; the dominance tests use frozen DINOv2 features, and the one fine-tuned candidate (a
paste-based model) was not a fair test because its pasting could not balance the traps.
